\chapter{Neue Drucksachen}
\quellseite{11}

Fünfzehnter Synodalbericht des Nebraska-Districts der deutschen ev.-luth. Synode von Missouri, Ohio u. a. Staaten. 1903. St. Louis, Mo. Concordia Publishing House. 120 Seiten 9×6. Preis: 25 Cts.

Dieser Bericht enthält ein sehr lehrreiches Referat Dr. Gräbners. An der Hand der Apostelgeschichte und der apostolischen Briefe wird uns das kirchliche Leben der apostolischen Zeit beschrieben. Da wird uns in lebendiger Weise vor Augen geführt, wie es damals stand mit der kirchlichen Ortsgemeinde, wie das kirchliche Amt verwaltet wurde. Wir hören ferner, wie die ersten Gemeinden ihren Gottesdienst eingerichtet hatten und zum Lob und Preis Gottes und zur Erbauung der Gemeinde hielten, wie Gottes Wort bei ihnen reichlich im Schwange ging und die heiligen Sacramente verwaltet wurden. Und endlich wird uns vorgeführt, wie auch die ersten Gemeinden unter Leitung der Apostel es lernten, christliche Kirchenzucht zu üben, um die sündigenden Brüder zu gewinnen oder die Unbußfertigen von sich auszuschließen. Dies alles wird in unserer Zeit und in unseren Gemeinden vorgehalten als ein Spiegelbild. Es wäre zu wünschen, daß dieser Bericht viele aufmerksame Leser fände. Auch der Missionsbericht ist interessant und gibt einen Einblick in die Thätigkeit des Districts. G. M.

Predigt über 5 Mos. 33, 24. 25., gehalten bei Gelegenheit des sechzigjährigen Jubiläums der ev.-luth. Dreieinigkeits-Gemeinde in New York am 7. Juni 1903 von einem Kinde der Gemeinde (P. P. Hirten). 7 Seiten 9½×6. Zu beziehen von Rev. O. Graesser, 602 E. 9th St., New York City. Preis: 5 Cts. portofrei.

Eine schöne, lesenswerthe Jubelpredigt, einer unserer alten Synodal-Gemeinden von einem ihrer dankbaren Kinder gehalten, die in dem Segenswunsch gipfelt: Dein Alter sei wie deine Jugend! Der Prediger zeigt an der Hand des sehr passenden Textes, was die Meinung dieses Wunsches sei und was Seitens der Gemeinde geschehen müsse, damit sich das Gewünschte an und bei ihr erfülle. L. F.

Predigt, gehalten am 25. October 1903 zum hundertjährigen Jubiläumsfest der Gründung der Stadt Chicago von Karl Schmidt, Pastor der ev.-luth. St. Jacobi-Gemeinde daselbst. 14 Seiten 7×5. Zu haben im Concordia Publishing House, St. Louis, Mo. Preis: 5 Cts.

Eine Gelegenheitsrede, die nicht in der bekannten Schematicmenie einherfährt, sondern den rechten Ton anschlägt. Auf Grund der Textesworte Ps. 143, 5.: „Ich gedenke an die vorigen Zeiten, ich rede von allen deinen Thaten, und sage von den Werken deiner Hände“, behandelt sie das Thema: „Unsere Theilnahme an der Jubelfeier unserer Stadt“, und zeigt, worüber und wie gejubelt werden soll. L. F.

The Lutheran Clarion. 4 Seiten 9×6. Zu beziehen von Rev. W. Hallerberg jun., 728 S. 8th St., Quincy, Ill. Preis: 50 Exemplare monatlich pro Jahr \$5.00.

Unter diesem Titel geben unsere Pastoren in Quincy und Umgegend ein englisches Blättchen heraus, das besonders zum fleißigen, regelmäßigen und andächtigen Besuch des Gotteshauses ermuntert und darum auch recht passend als Tractat an den Kirchstüren vertheilt werden kann, ähnlich wie der bekannte in Chicago erscheinende „Stadtmissionar“. Das Blatt ist namentlich da zu gebrauchen, wo mehr oder weniger regelmäßig englische Gottesdienste gehalten werden, die von Fremden besucht werden. In dem oben angegebenen Preis sind die für jede Gemeinde besonders gedruckten Bekanntmachungen des Ortes und der Zeit der Gottesdienste mit eingeschlossen. L. F.

Vom Schriftenverein der sep. ev.-luth. Gemeinden in Sachsen aus Zwickau sind die beiden folgenden Schriften eingegangen, die bei unserem Concordia Publishing House, St. Louis, Mo., bestellt werden können:

Warum hangen wir so fest an der lutherischen Kirche? Beantwortet von C. F. W. Walther, weiland Professor der Theologie in St. Louis, Mo. Vierte Auflage. 20 Seiten 9×5½. Preis: 10 Cts.

Dies ist ein schon wiederholt aufgelegter ganz trefflicher Tractat unseres unvergesslichen Dr. Walther, der überall nur mit großem Segen gelesen werden kann und in jedem Leser neue Liebe zu seiner Kirche erwecken wird. Walther beantwortet die gestellte Frage namentlich mit zwei weitläufig ausgeführten Gründen, nämlich erstens, weil die lutherische Kirche sich in ihrer Lehre ganz genau und ganz allein nach Gottes Wort richtet, und zweitens, weil sie das Evangelium im engeren Sinne, die Lehre von der Rechtfertigung allein durch den Glauben, so rein und lauter, so fröhlich und kräftig, so reich und in ihrer ganzen Fülle vorträgt. Am Schlusse gibt Walther noch kurz 15 weitere Gründe an. Das Schriftchen sei wieder bestens empfohlen.

Wie wird man vor Gott gerecht und selig? Synodalpredigt über das Evangelium am 11. Sonntag nach Trinitatis von R. H. Beer, ev.-luth. Pastor in Lincoln Township, Plymouth Co., Iowa. 16 Seiten 7×5. Preis: 5 Cts.

Diese Predigt wurde im Jahre 1900 bei Gelegenheit der Synodalversammlung unsers Iowa-Districts gehalten. Sie behandelt auf Grund des Evangeliums vom Pharisäer und Zöllner in lehrhafter Weise den Grundartikel von der Rechtfertigung, indem sie zeigt, wie der Mensch, der gerecht und selig werden will, in rechter Erkenntniß seiner Sünde und in rechter Erkenntniß der Gnade Gottes stehen muß. L. F.

Geistliche Lieder für Männerchöre. 4. Heft. Concordia Publishing House. St. Louis, Mo. 12 Seiten 11×8. Preis: 20 Cts.; Dutzendpreis: \$1.50 und Porto.

Dies neueste Heft der schon wiederholt warm empfohlenen Sammlung bringt eine ganz vorzügliche Auswahl von Passions- und Confirmationsliedern. Die einzelnen Nummern sind: Siehe, das ist Gottes Lamm (C. Chr. Weeber); Die sieben Worte (R. Palme); Ehre sei dir, Christe (H. Wehe); Herr Jesu, du lieb ich (C. Palmer); Stärkt uns, Mittler (Knecht); Verlaß mich nicht (S. Zepp); Wie könnt ich sein vergessen (Orlando di Lasso). Wir lenken insbesondere die Aufmerksamkeit auf den ergreifenden Charfreitagsgesang: Die sieben Worte. Das Stück wird von zwei Chören gesungen. Der erste Chor singt die sieben Worte; zwischen den einzelnen Worten singt dann der zweite Chor, der durch ein Doppelquartett vertreten wird, dreimal: „Herr (Christe), erbarme dich unser“ und die drei Verse von „Christe, du Lamm Gottes“. Sehr geschickt ist dies zusammengebogen. Zum Schlusse, nach dem siebenten Wort: „Vater, in deine Hände befehle ich meinen Geist“, singen beide Chöre zusammen den Choral: „Wann ich einmal soll scheiden“ in dem bekannten eigenartig schönen Tonsatz aus der „Matthäuspassion“ von J. S. Bach. Gut vorgetragen, wird das Stück tiefen Eindruck machen. L. F.
